\chapter[Terminal, naturality, and limit]{Multivalued terminal, unique naturality and passage to the limit}
\label{chap:pdfv2-terminal-naturalidad-limite}

\begin{thesisbox}
The discrete continuum produces a single system of histories, extensions, memories, and terminals. The internal discriminants \(\mathrm{sp},\mathrm{sol},\mathrm{gau},\mathrm{coh},\mathrm{det}\) name five characteristics that are simultaneous within that same system: they neither fix distinct inputs nor authorize the selection of different histories. In this chapter, the terminal preserves all continuations, the connection remains a nonadic correspondence, and naturality is expressed by a single identity.
\end{thesisbox}

\section{Surviving histories with extension witnesses}

Let
\(\widehat x_k\in\widehat{\mathcal X}^{\rm tot}_k\) be the total state of
\eqref{eq:continuo-comun-dominio-total}. For \(N\geq k\), we define without
any other predicate of membership
\begin{equation}
\label{eq:pdfv2-terminal-historias-finita}
 \mathscr H_{k,N}(\widehat x_k)
 :=\operatorname{Term}_{k,N}(\widehat x_k).
\end{equation}
Each
\(h=(\varepsilon_{k+1},\ldots,\varepsilon_N)\) on the right contains through
\eqref{eq:continuo-comun-vertice-arbol} its intermediate states, the witnesses
of \(\operatorname{Ext}\), the carries, boundaries, routes, and ledger.
Each edge satisfies the phase action
\eqref{eq:continuo-comun-accion-fase-tpk}; hence every nine-step subpath
determines, with its nine emissions and not merely with its length, an
element of \(\mathsf Z^\Gamma_{9,j}\) and a witness of
\(\Gamma_{9,j}\) through
\eqref{eq:continuo-comun-gamma9-apice}--
\eqref{eq:continuo-comun-gamma9-relacion}. Therefore, Ext and nonadic holonomy
arise from the same phase-compatible word. To \eqref{eq:pdfv2-terminal-historias-finita} there is adjoined neither a survival filter nor a
second condition \(\Gamma_9\).

For \(k\leq\ell\leq N\), let
\begin{equation}
\label{eq:pdfv2-terminal-prefijos-intermedios}
 \mathscr P_{\ell\mid k}(\widehat x_k)
 :=\{h_{[k,\ell]}:h\in\mathscr H_{k,N}(\widehat x_k)
       \text{ for some }N\geq\ell\}.
\end{equation}
For \(p\in\mathscr P_{\ell\mid k}(\widehat x_k)\), we use the abbreviation
\begin{equation}
\label{eq:pdfv2-terminal-continuaciones-prefijo}
 \mathscr H_{\ell,N}(p)
 :=\{q:p\star q\in\mathscr H_{k,N}(\widehat x_k)\}.
\end{equation}
The prolongation neutral of \eqref{eq:continuo-comun-seccion-total} extends
each prefix to cualquier horizon. Consequently,
\begin{equation}
\label{eq:pdfv2-terminal-no-vacuidad-historias}
 \mathscr H_{k,N}(\widehat x_k)\neq\varnothing,
 \qquad
 \mathscr P_{\ell\mid k}(\widehat x_k)\neq\varnothing,
\end{equation}
by \cref{prop:continuo-comun-arbol-finito-no-vacio}.

The truncation
\begin{equation}
\label{eq:pdfv2-terminal-truncamiento-historias}
 \pi_{N+1,N}^{\mathscr H}:
 \mathscr H_{k,N+1}(\widehat x_k)
 \longrightarrow\mathscr H_{k,N}(\widehat x_k)
\end{equation}
is exactly \eqref{eq:continuo-comun-proyeccion-terminal} and is
surjective by \cref{lem:continuo-comun-terminal-sobreyectivo}. Survival is, by
\eqref{eq:continuo-comun-surv-igual-total}, a demonstrated output of the total
domain and not a prior selection.

\section{Generated terminal record and its \(13\) coordinates}

Let us fix a history
\[
 h=(\varepsilon_{j+1}:x_j\longrightarrow x_{j+1})_{j=k}^{N-1},
 \qquad
 \gamma_{k,k}:=I_{x_k},\qquad
 \gamma_{k,j+1}:=\gamma_{k,j}\star\varepsilon_{j+1},
\]
The operator notation
\(U_{\gamma_{k,j+1}}
=U_{\varepsilon_{j+1}}U_{\gamma_{k,j}}\) preserves the order of application,
whereas \(\star\) preserves the chronological order of the word. We further fix
a coefficient level \(r\geq1\). For brevity, we write
\[
 C_{x_j,r}:=C_{j,r}(x_j),\qquad
 F_{\varepsilon_{j+1}}:C_{x_{j+1},r}\longrightarrow C_{x_j,r}
\]
for the complex and the map of complexes of
\eqref{eq:pdfv2-cor-cono-celda} and
\eqref{eq:pdfv2-cor-refinamiento-celda}. The type of the level record
\(j\), denoted \(\mathfrak R_{j,r}\), is a dependent record with the
following thirteen fields. In each field we indicate the type, the seed and the
action of an edge \(\varepsilon:=\varepsilon_{j+1}\).

The state \(x_k\) already contains the complete word
\(\gamma_{0,k}=\varepsilon_1\star\cdots\star\varepsilon_k\). Therefore a
``seed in \(k\)'' never restarts the genealogy: it is the fold of the same
updates from the root to \(k\). We denote
\begin{equation}
\label{eq:pdfv2-terminal-prefijo-heredado}
 \begin{gathered}
 M_k^\pm:=\sum_{i=1}^{k}b_{\varepsilon_i}^\pm,qquad
 D_k:=\sum_{i=1}^{k}b_{\varepsilon_i},qquad
 H_k:=\sum_{i=1}^{k}|b_{\varepsilon_i}|,\\
 \mathbf c_{\leq k}:=
  \bigl(c_r(\varepsilon_i,\gamma_{0,i-1})\bigr)_{1\leq i\leq k},
 \qquad
 \mathbf k_{\leq k}:=
  \bigl(\kappa_{i-1}^{\rm can}\bigr)_{1\leq i\leq k},\\
 \mathbf F_{\leq k}:=
  \bigl(F_{\varepsilon_i}:C_{x_i,r}\to C_{x_{i-1},r}\bigr)_{1\leq i\leq k}.
 \end{gathered}
\end{equation}
All empty lists appearing below are understood literally only
in \(k=0\); for \(k>0\) they are replaced by prefixes inherited from
\eqref{eq:pdfv2-terminal-prefijo-heredado} and from the ledger of
\(\gamma_{0,k}\). Thus the record does not preserve an opaque copy of \(x_k\),
but neither does it lose its APP--TRIT--TPK past.

\paragraph{Fiber \(\mathsf F\): local domain and conservation under prolongation}
The elementary fiber is
\begin{equation}
\label{eq:pdfv2-terminal-tipo-fibra}
 \operatorname{Fib}_r(x):=
 \bigl(\mathsf H_x^\Sigma\oplus\mathsf H_x^\Pi,\,
       \mathsf M_x,\,
       \operatorname{Cell}_{\rm TPK}(x),\,C_{x,r}\bigr).
\end{equation}
The field \(\mathsf F_{k,j}\) is a word in the category of these fibers.
Its seed and update are
\begin{align}
 \mathsf F_{k,k}
 &:=\bigl(\operatorname{Fib}_r(x_0),I\bigr)\star
 \mathop{\star}_{i=1}^{k}
 \bigl(\operatorname{Fib}_r(x_i),U_{\varepsilon_i},
 \mathsf C_{\varepsilon_i},\operatorname{Cell}_{\rm TPK}(\varepsilon_i),
 F_{\varepsilon_i}\bigr),
 \label{eq:pdfv2-terminal-semilla-f}\\
 \operatorname{Upd}^{\mathsf F}_{\varepsilon}\mathsf F_{k,j}
 &:=
 \mathsf F_{k,j}\star
 \bigl(\operatorname{Fib}_r(x_{j+1}),U_\varepsilon,
 \mathsf C_\varepsilon,\operatorname{Cell}_{\rm TPK}(\varepsilon),
 F_\varepsilon\bigr).
 \label{eq:pdfv2-terminal-update-f}
\end{align}
It is nonempty because the two sheets, the memory, the cellular unit
\eqref{eq:pdfv2-cor-cell-unidad}, and the port complex have already been
constructed. The composition of its five maps proves naturality.

\paragraph{Extensions \(\mathsf E\): morphisms of compatible prolongation of the state}
The codomain is the category of typed relation words. We set
\begin{align}
 \mathsf E_{k,k}
 &:=(I_{x_0},\operatorname{Ext}_{I_{x_0}})\star
 \mathop{\star}_{i=1}^{k}
 \bigl(\varepsilon_i,\operatorname{Ext}_{\varepsilon_i},
 \operatorname{Ext}_{\varepsilon_i}\circ
 \operatorname{Ext}_{\gamma_{0,i-1}}
 \subseteq\operatorname{Ext}_{\gamma_{0,i}}\bigr),
 \label{eq:pdfv2-terminal-semilla-e}\\
 \operatorname{Upd}^{\mathsf E}_{\varepsilon}\mathsf E_{k,j}
 &:=
 \mathsf E_{k,j}\star
 \bigl(\varepsilon,\operatorname{Ext}_{\varepsilon},
 \operatorname{Ext}_{\varepsilon}\circ
 \operatorname{Ext}_{\gamma_{0,j}}
 \subseteq\operatorname{Ext}_{\gamma_{0,j+1}}\bigr).
 \label{eq:pdfv2-terminal-update-e}
\end{align}
Identity and neutral emission prove nonemptiness; associativity of the
composition relational proof naturality.

\paragraph{Relations \(\mathsf R\): incidence, orientation, and compatibility of composition}
Let \(\operatorname{RelBlock}_{j,r}\) be the type of blocks of equations
labeled by the edge and the prefix over which they are verified, and let
\(\operatorname{Pres}_{j,r}:=\operatorname{List}
(\operatorname{RelBlock}_{j,r})\). The unit block and the seed are
\begin{equation}
\label{eq:pdfv2-terminal-semilla-r}
 \mathsf B_I:=(I_{x_0};U_I=I,\ \mathsf U_I=I,\ \kappa(I)=0,\ \partial I=0),
 \qquad
 \mathsf R_{k,k}:=(\mathsf B_I)\star
 \mathop{\star}_{i=1}^{k}
 \mathsf B_R(\varepsilon_i;\gamma_{0,i-1}).
\end{equation}
For an edge \(\varepsilon:x_j\to x_{j+1}\), the new block is the
labeled tuple
\begin{equation}
\label{eq:pdfv2-terminal-bloque-r}
 \begin{split}
 \mathsf B_R(\varepsilon;\gamma_{0,j}):=\bigl(
 &\varepsilon,\gamma_{0,j};\ 
 \delta^\diamond(\varepsilon)
 =\operatorname{res}^\diamond(\varepsilon)
  +9\operatorname{quo}^\diamond(\varepsilon),\\
 &\Delta(\varepsilon)=\operatorname{or}(\varepsilon)
  +3\kappa(\varepsilon),\quad
 U_{\gamma_{0,j+1}}=U_\varepsilon U_{\gamma_{0,j}},\\
 &\mathsf U_{\gamma_{0,j+1}}
 =\mathsf U_\varepsilon\mathsf U_{\gamma_{0,j}},\quad
 \kappa(\gamma_{0,j+1})
 =\kappa(\varepsilon)+\mathsf C_\varepsilon\kappa(\gamma_{0,j}),\\
 &\partial\gamma_{0,j+1}
 =\partial\varepsilon+\partial\gamma_{0,j},\quad
 \operatorname{Cpl}_\varepsilon
 \text{ of }\eqref{eq:pdfv2-cor-acoplamientos-unicos}
 \bigr).
 \end{split}
\end{equation}
The update concatenates the block without losing date or multiplicity:
\begin{equation}
\label{eq:pdfv2-terminal-update-r}
 \operatorname{Upd}^{\mathsf R}_{\varepsilon}\mathsf R_{k,j}
 :=\mathsf R_{k,j}\star\mathsf B_R(\varepsilon;\gamma_{0,j}).
\end{equation}
The truncation \(\operatorname{tr}^{\mathsf R}_{j+1,j}\) removes the final
labeled block. Moreover,
\(a_*\mathsf B_R(\varepsilon;\gamma)
=\mathsf B_R(a\varepsilon;a\gamma)\), because each equation is natural over
the same edge and the same prefix.

\paragraph{Numbers \(\mathsf N\): publication of value with conservation of genealogy}
For an edge \(\varepsilon:x_j\to x_{j+1}\), we define the numerical datum
\begin{equation}
\label{eq:pdfv2-terminal-numero-arista}
 \nu(\varepsilon):=
 \bigl(j+1,\mathsf Q_{j+1},g_{j+1},q^{(9)}_{j+1},\beta_{j+1},
 \operatorname{res}^{\Sigma},\operatorname{quo}^{\Sigma},
 \operatorname{res}^{\Pi},\operatorname{quo}^{\Pi},
 \operatorname{or},\kappa,\Delta\bigr)(\varepsilon).
\end{equation}
Its type is
\[
 \mathbb N\times\widetilde{\mathcal O}_9\times C_9\times\mathbb N
 \times\mathbb Z/12\mathbb Z
 \times\{0,\ldots,8\}^2\times\mathbb Z^5.
\]
The seed contains the complete numerical list of the prefix \(0\to k\);
the update adds exactly the datum of the next edge:
\begin{equation}
\label{eq:pdfv2-terminal-update-n}
 \mathsf N_{k,k}:=
 ((0,\mathsf Q_0,g_0,q^{(9)}_0,\beta_0;0,\ldots,0))
 \star\mathop{\star}_{i=1}^{k}\nu(\varepsilon_i),\qquad
 \operatorname{Upd}^{\mathsf N}_{\varepsilon}\mathsf N_{k,j}
 :=\mathsf N_{k,j}\star\nu(\varepsilon).
\end{equation}
Totality proceeds from the divisions APP, the division TRIT, and the odometer TPK.

\paragraph{Orientation \(\mathsf O\): sign transport and leaf selection}
The type of an orientation state is
\[
 \mathbb Z\times
 (\mathsf H^\Sigma_x\oplus\mathsf H^\Pi_x)
 \times\mathbb Z[\mathsf A],
\]
with seed and update
\begin{align}
 \mathsf O_{k,k}&:=
 \left(\sum_{i=1}^{k}\operatorname{or}(\varepsilon_i),
       \operatorname{Sh}(x_k),\partial\gamma_{0,k}\right),
 \label{eq:pdfv2-terminal-semilla-o}\\
 \operatorname{Upd}^{\mathsf O}_{\varepsilon}
 (\mathfrak o,\operatorname{Sh},\mathfrak b)
 &:=
 \bigl(\mathfrak o+\operatorname{or}(\varepsilon),
 U_\varepsilon\operatorname{Sh},
 \mathfrak b+\partial\varepsilon\bigr).
 \label{eq:pdfv2-terminal-update-o}
\end{align}
The APP leaves prove nonemptiness. Functoriality of \(U\), additivity of
\(\operatorname{or}\), and telescoping of the boundary prove naturality.

\paragraph{Carry \(\mathsf K\): update of the quotient and residual memory}
We use the typed integer lifts in
\eqref{eq:pdfv2-cor-carry-matricial-def}. The field has type
\[
 \mathsf M_{x_j}\times
 \operatorname{Hom}(\mathsf M_{x_0},\mathsf M_{x_j})
 \times\operatorname{Mat}_5(\mathbb Z)
 \times\operatorname{List}(\operatorname{Mat}_5(\mathbb Z))
 \times\operatorname{List}(\mathbb Q_{\geq0}).
\]
For \(\gamma_{k,j}=\varepsilon_j\cdots\varepsilon_{k+1}\), we explicitly
fix
\begin{equation}
\label{eq:pdfv2-terminal-lift-entero-ruta}
 \widetilde L_{\gamma_{k,k}}:=I_5,\qquad
 \widetilde L_{\gamma_{k,j}}
 :=\widetilde L_r(\varepsilon_j)\cdots
   \widetilde L_r(\varepsilon_{k+1}).
\end{equation}
For the inherited prefix we also fix
\[
 \widetilde L_{\gamma_{0,k}}
 :=\widetilde L_r(\varepsilon_k)\cdots
   \widetilde L_r(\varepsilon_1).
\]
The seed and update are
\begin{align}
 \mathsf K_{k,k}&:=
 \bigl(\kappa(\gamma_{0,k}),\mathsf C_{\gamma_{0,k}},
       \widetilde L_{\gamma_{0,k}},
       \mathbf c_{\leq k},\mathbf k_{\leq k}\bigr),
 \label{eq:pdfv2-terminal-semilla-k}\\
 \operatorname{Upd}^{\mathsf K}_{\varepsilon}
 (\kappa_\gamma,\mathsf C_\gamma,\widetilde L_\gamma,
   \mathbf c,\mathbf k)
 &:=
 \bigl(\kappa(\varepsilon)+\mathsf C_\varepsilon\kappa_\gamma,\,
 \mathsf C_\varepsilon\mathsf C_\gamma,\,
 \widetilde L_r(\varepsilon)\widetilde L_\gamma,\,
 \mathbf c\star c_r(\varepsilon,\gamma),\,
 \mathbf k\star\kappa_j^{\rm can}\bigr).
 \label{eq:pdfv2-terminal-update-k}
\end{align}
The identity provides the seed; the two cocycles
\eqref{eq:continuo-comun-cociclo-carry} and
\eqref{eq:pdfv2-cor-cociclo-matricial} prove independence of
parenthesization and naturality.

\paragraph{Memory \(\mathsf M\): record of the traversal and compatibility under truncation}
Its type at the same time preserves affine memory, ledger, and signed memories:
\begin{equation}
\label{eq:pdfv2-terminal-tipo-ledger-entry}
 \operatorname{LedgerEntry}:=
 \{\lambda(\varepsilon):\varepsilon\text{ typed elementary emission}\}.
\end{equation}
\[
 \mathsf M_x\times\operatorname{List}(\operatorname{LedgerEntry})
 \times\mathbb N^2\times\mathbb Z\times\mathbb N.
\]
We set
\begin{align}
 \mathsf M_{k,k}&:=
 \bigl(m_k,\operatorname{Led}(\gamma_{0,k}),
       M_k^+,M_k^-,D_k,H_k\bigr),
 \label{eq:pdfv2-terminal-semilla-m}\\
 \operatorname{Upd}^{\mathsf M}_{\varepsilon}
 (m,\mathbf{Led},M^+,M^-,D,H)
 &:=
 \bigl(\mathsf C_\varepsilon m+\kappa(\varepsilon),\,
 \mathbf{Led}\star\lambda(\varepsilon),\,
 M^++b_\varepsilon^+,\,M^-+b_\varepsilon^-,\,
 D+b_\varepsilon,\,H+|b_\varepsilon|\bigr).
 \label{eq:pdfv2-terminal-update-m}
\end{align}
The affine action and \eqref{eq:pdfv2-cor-memorias} prove totality and
naturality; the past is not recomputed from the last value.

\paragraph{Boundary \(\mathsf B\): gluing data and survival of the extension}
The type is a boundary chain accompanied by a word of complejos
based and chain maps. Its seed and update are
\[
 \mathsf B_{k,j}\in
 \operatorname{Rec}[\mathfrak b\in\mathbb Z[\mathsf A_j];
 C_{x_j,r}\in\operatorname{Ch}_{\rm bas};
 \mathbf F\in\operatorname{Path}_{\operatorname{Ch}_{\rm bas}}
 (C_{x_j,r},C_{x_0,r})].
\]
\begin{align}
 \mathsf B_{k,k}&:=
 \bigl(\partial\gamma_{0,k},C_{x_k,r},\mathbf F_{\leq k}\bigr),
 \label{eq:pdfv2-terminal-semilla-b}\\
 \operatorname{Upd}^{\mathsf B}_{\varepsilon}
 (\mathfrak b,C_{x_j,r},\mathbf F)
 &:=
 \bigl(\mathfrak b+\partial\varepsilon,\,
 C_{x_{j+1},r},\,\mathbf F\star F_\varepsilon\bigr).
 \label{eq:pdfv2-terminal-update-b}
\end{align}
Non-emptiness proceeds from the common ports; \(\partial^2=0\),
\(F_{\beta\varepsilon}=F_\varepsilon F_\beta\), and
\eqref{eq:continuo-comun-frontera-telescopica} prove compatibility.

\paragraph{Route \(\mathsf G\): observable history of the surviving extension}
The field is the list of prefixes in the path category:
\begin{equation}
\label{eq:pdfv2-terminal-update-g}
 \mathsf G_{k,k}:=(\gamma_{0,i})_{0\leq i\leq k},\qquad
 \operatorname{Upd}^{\mathsf G}_{\varepsilon}\mathsf G_{k,j}
 :=\mathsf G_{k,j}\star\gamma_{0,j+1}.
\end{equation}
Identity, associative concatenation, and
\(a(\varepsilon\gamma)=a(\varepsilon)a(\gamma)\) prove the three
required properties.

\paragraph{Multisection \(\mathsf X\): stable finite fiber and obstruction to pointwise selection}
For \(M\geq j\), we define by action
\begin{equation}
\label{eq:pdfv2-terminal-multiseccion-interna}
 \mathsf X_{k,j}(M):=
 \{\gamma_{k,j}\star q:
 q\in\operatorname{Term}_{j,M}(x_j)\}.
\end{equation}
Thus
\[
 \mathsf X_{k,k}(M)=\operatorname{Term}_{k,M}(x_k),
\]
and
\begin{equation}
\label{eq:pdfv2-terminal-update-x}
 \operatorname{Upd}^{\mathsf X}_{\varepsilon}\mathsf X_{k,j}(M)
 :=
 \{\gamma_{k,j}\star\varepsilon\star q:
 q\in\operatorname{Term}_{j+1,M}(x_{j+1})\}.
\end{equation}
The neutral prolongation proves nonemptiness. Deleting the final step and acting
by a symmetry commute literally with concatenation.

\paragraph{Survival \(\mathsf S\): persistence under refinement and passage to the limit}
Here, a proposition is not stored, but witnesses. Let
\begin{equation}
\label{eq:pdfv2-terminal-testigo-neutro}
 s_{j,j}:=I_{x_j},\qquad
 s_{j,M+1}:=s_{j,M}\star
 \varepsilon_M^0(t(s_{j,M})).
\end{equation}
The seed is
\(\mathsf S_{k,k}:=(s_{k,M})_{M\geq k}\), and
\begin{equation}
\label{eq:pdfv2-terminal-update-s}
 \operatorname{Upd}^{\mathsf S}_{\varepsilon}\mathsf S_{k,j}
 :=(\gamma_{k,j}\star\varepsilon\star s_{j+1,M})_{M\geq j+1}.
\end{equation}
By construction
\(\pi_{M+1,M}s_{j,M+1}=s_{j,M}\); the naturality of neutral emission
proves equivariance.

\paragraph{Terminal \(\mathsf T\): universality of the limit object and factorization of readers}
The terminal is the inverse system, not a selected queue:
\begin{equation}
\label{eq:pdfv2-terminal-sistema-inverso-campo}
 \mathsf T(x_j):=
 \bigl((\operatorname{Term}_{j,M}(x_j))_{M\geq j},
       (\pi_{M+1,M})_{M\geq j}\bigr).
\end{equation}
For each \(M\geq j\), the map of prefix tipado is
\begin{equation}
\label{eq:pdfv2-terminal-update-t}
 \operatorname{pref}_{\gamma_{k,j},M}:
 \operatorname{Term}_{j,M}(x_j)\longrightarrow
 \operatorname{Term}_{k,M}(x_k),\qquad
 q\longmapsto\gamma_{k,j}\star q.
\end{equation}
We define the field state and its update by
\begin{align}
 \mathsf T_{k,j}
 &:=
 \left(\mathsf T(x_j),
   (\operatorname{pref}_{\gamma_{k,j},M})_{M\geq j}\right),
 \label{eq:pdfv2-terminal-estado-t}\\
 \operatorname{Upd}^{\mathsf T}_{\varepsilon}\mathsf T_{k,j}
 &:=
 \left(\mathsf T(x_{j+1}),
   (\operatorname{pref}_{\gamma_{k,j+1},M})_{M\geq j+1}\right),
 \qquad \gamma_{k,j+1}:=\gamma_{k,j}\star\varepsilon.
 \label{eq:pdfv2-terminal-update-t-explicita}
\end{align}
In particular, \(\mathsf T_{k,k}\) uses
\(\gamma_{k,k}=I_{x_k}\). The witnesses
\eqref{eq:pdfv2-terminal-testigo-neutro} and the surjectivity of \(\pi\)
prove that the system is nonempty; level by level,
\begin{equation}
\label{eq:pdfv2-terminal-prefijo-natural}
 \pi^{\mathscr H}_{M+1,M}
 \operatorname{pref}_{\gamma_{k,j},M+1}
 =\operatorname{pref}_{\gamma_{k,j},M}
  \pi^{\mathscr H}_{M+1,M},
\end{equation}
because both traversals only remove the last emission of \(q\).

\paragraph{13. Terminal reader \(\mathsf L\): evaluation of the compatible extension and conservation of genealogy}
The reader contains no classical object. Its type is the category of
paths of the unique dependent constructor
\(\mathfrak D_{j,r}\) of \eqref{eq:pdfv2-cor-constructor-unico}. We set
\begin{equation}
\label{eq:pdfv2-terminal-transicion-lector-tipado}
 u^{\mathfrak D}_{j+1,j}
 :=u_{(j+1,r),(j,r)}:
 \mathfrak D_{j+1,r}(x_{j+1})
 \longrightarrow\mathfrak D_{j,r}(x_j),
\end{equation}
the component of the common transport
\eqref{eq:pdfv2-cor-naturalidad-unica} that truncates the last emission, reduces
the coefficients and precomposes the ports. By composition,
\begin{equation}
\label{eq:pdfv2-terminal-transicion-lector-composicion}
 u^{\mathfrak D}_{m,k}
 =u^{\mathfrak D}_{\ell,k}u^{\mathfrak D}_{m,\ell}
 \qquad(m\geq\ell\geq k).
\end{equation}
\begin{align}
 \mathsf L_{k,k}
 &:=(\mathfrak D_{0,r}(x_0))\star
 \mathop{\star}_{i=1}^{k}
 \bigl(\mathfrak D_{i,r}(x_i),u^{\mathfrak D}_{i,i-1},
       \mathcal O(\mathfrak e_{\varepsilon_i})\bigr),
 \label{eq:pdfv2-terminal-semilla-l}\\
 \operatorname{Upd}^{\mathsf L}_{\varepsilon}\mathsf L_{k,j}
 &:=
 \mathsf L_{k,j}\star
 \bigl(\mathfrak D_{j+1,r}(x_{j+1}),
 u^{\mathfrak D}_{j+1,j},\mathcal O(\mathfrak e_\varepsilon)\bigr).
 \label{eq:pdfv2-terminal-update-l}
\end{align}
The totality of \(\mathfrak D\) and
\eqref{eq:pdfv2-terminal-transicion-lector-composicion} proves non-emptiness
and naturality.

The thirteen preceding seeds form a single element.
\begin{equation}
\label{eq:pdfv2-terminal-semilla-trece}
 \mathfrak r^0_{k,r}:=
 \langle
 \mathsf F_{k,k};\mathsf E_{k,k};\mathsf R_{k,k};\mathsf N_{k,k};
 \mathsf O_{k,k};\mathsf K_{k,k};\mathsf M_{k,k};\mathsf B_{k,k};
 \mathsf G_{k,k};\mathsf X_{k,k};\mathsf S_{k,k};\mathsf T_{k,k};
 \mathsf L_{k,k}
 \rangle\in\mathfrak R_{k,r}.
\end{equation}
The action of an edge is the total application.
\begin{equation}
\label{eq:pdfv2-terminal-actualizacion-unica}
 \operatorname{Upd}_\varepsilon:
 \mathfrak R_{j,r}(\gamma_{0,j})
 \longrightarrow\mathfrak R_{j+1,r}(\gamma_{0,j+1}),\qquad
 \operatorname{Upd}_\varepsilon:=
 \langle\operatorname{Upd}^{\mathsf F}_\varepsilon,\ldots,
 \operatorname{Upd}^{\mathsf L}_\varepsilon\rangle,
\end{equation}
with the thirteen actions
\eqref{eq:pdfv2-terminal-update-f}--
\eqref{eq:pdfv2-terminal-update-l}. Finally, the terminal register is
\emph{generated} by the fold
\begin{equation}
\label{eq:pdfv2-terminal-registro-trece}
 \boxed{
 \mathfrak r_{k,N}(h):=
 \operatorname{Upd}_{\varepsilon_N}\circ\cdots\circ
 \operatorname{Upd}_{\varepsilon_{k+1}}(\mathfrak r^0_{k,r}).}
\end{equation}
It is not the unpacking of an intact copy of the state.
On the image of this fold, the truncation is defined componentwise
\begin{equation}
\label{eq:pdfv2-terminal-truncamiento-trece-def}
 \operatorname{tr}_{13;j+1,j}
 \mathfrak r_{k,j+1}(h\star\varepsilon)
 :=\mathfrak r_{k,j}(h).
\end{equation}
It is well defined: the fields \(\mathsf E\) and \(\mathsf G\) preserve the
ordered list of emissions and prefixes, so that equality of two
records entails equality of their labeled histories. In
\(\mathsf F,\mathsf E,\mathsf R,\mathsf N,\mathsf G,\mathsf T,\mathsf L\),
the last entry or block is deleted; in \(\mathsf O,\mathsf K,\mathsf M,
\mathsf B\), the cumulative formula is applied to the recovered prefix; and in
\(\mathsf X,\mathsf S\), the cone of tails is restricted. Neither an
idempotent union is inverted nor a continuation selected.

\begin{proposition}[Totality and naturalness of the thirteen updates]
\label{prop:pdfv2-terminal-trece-total-natural}
All the seeds of \eqref{eq:pdfv2-terminal-semilla-trece} exist, each
\(\operatorname{Upd}_\varepsilon\) is total, and, for every coherent symmetry
\(a\) and every horizon truncation,
\begin{equation}
\label{eq:pdfv2-terminal-trece-naturalidad}
 a_*\operatorname{Upd}_\varepsilon
 =\operatorname{Upd}_{a\varepsilon}a_*,
 \qquad
 \operatorname{tr}_{13;j+1,j}
 \operatorname{Upd}_\varepsilon\mathfrak r_{k,j}(h)
 =\mathfrak r_{k,j}(h).
\end{equation}
\end{proposition}

\begin{proof}
The existence of each seed was exhibited in the thirteen paragraphs. In
\(\mathsf F,\mathsf E,\mathsf G,\mathsf X,\mathsf S,\mathsf T,\mathsf L\),
the formula is concatenation with an already typed map. In
\(\mathsf R,\mathsf N,\mathsf O,\mathsf K,\mathsf M,\mathsf B\), the
formula is, respectively, concatenation of a labeled block of equations,
adjunction of a numerical datum, transport of sheets, composition
of cocycles, affine action, and chain map. The naturality identities
cited after each formula prove the first equality of
\eqref{eq:pdfv2-terminal-trece-naturalidad} componentwise. The
second is exactly
\eqref{eq:pdfv2-terminal-truncamiento-trece-def}; its sound definition uses the
ordered history preserved by \(\mathsf E\) and \(\mathsf G\).
\end{proof}

\begin{proposition}[Joint Certificate of the Five Internal Developments]
\label{prop:pdfv2-terminal-cinco-desarrollos-integrales}
The updates of the thirteen coordinates jointly transport:
the spectral--polar telescoping and return without reinitialization of
\cref{thm:continuo-sp-telescopia-no-autonoma,eq:continuo-sp-no-reset};
the orthogonal balance, the remark column, and the Stein identity of
\cref{eq:pdfv2-sol-ortogonalidad,eq:pdfv2-sol-stein}; direct
and adjoint naturality and the cellular retraction of
\cref{eq:gauint-naturalidad-directa-adjunta,eq:gauint-sdr}; totalization,
the cone of the nonadic return, and the complete fibers of
\cref{eq:pdfv2-coh-total-square-zero,eq:pdfv2-coh-terminal-cone,eq:pdfv2-coh-extension-multisection}; and the Alexander--Whitney
diagonal, Smith reduction, and the Bockstein tower of
\cref{eq:detint-aw-natural,eq:detint-smith,eq:detint-bock-leading}.
All these actions have as their domain the same register generated by
\eqref{eq:pdfv2-terminal-registro-trece}; none introduces an initial
subdomain, a privileged history, or a selection motivated by a subsequent
application.
\end{proposition}

\begin{proof}
Each cited identity acts on a coordinate of the common register, and its truncation maps are restrictions of \eqref{eq:pdfv2-terminal-truncamiento-trece-def}. The componentwise naturality proved in \eqref{eq:pdfv2-terminal-trece-naturalidad} permits their assembly into a single update. Equality of domain and preservation of the ordered history preclude replacing correlative generation with five retrospective projections.
\end{proof}

\section{Multivalued terminal without selector}

The finite terminal generated by \(\widehat x_k\) is the multisection
\begin{equation}
\label{eq:pdfv2-terminal-multivaluado}
 \mathfrak T_{k,N}(\widehat x_k)
 :=\left\{\mathfrak r_{k,N}(h):
            h\in\mathscr H_{k,N}(\widehat x_k)\right\}.
\end{equation}
The expression of the right returns the set whole. No contains a
rule that distinguishes a history privileged.

The terminal level environment is obtained by gathering those multisects without
identifying records:
\begin{equation}
\label{eq:pdfv2-terminal-ambiente-nivel}
 \mathfrak T_{k,N}
 :=\bigcup_{\widehat x_k\in\widehat{\mathcal X}^{\rm tot}_k}
   \mathfrak T_{k,N}(\widehat x_k).
\end{equation}

The linear version, when contributions must be added, is defined on
the free module by
\begin{equation}
\label{eq:pdfv2-terminal-evaluacion-todas}
 \begin{aligned}
 \mathsf E_{k,N}^{\rm all}:
 \mathbb Z[\widehat{\mathcal X}^{\rm tot}_k]&\longrightarrow
 \mathbb Z[\mathfrak T_{k,N}],\\
 [\widehat x_k]&\longmapsto
 \sum_{h\in\mathscr H_{k,N}(\widehat x_k)}
 [\mathfrak r_{k,N}(h)].
 \end{aligned}
\end{equation}
Each witness appears with its multiplicity. If exists a measure projective
positive \(\mu\), the version isometric uses all the probabilities
conditional:
\begin{equation}
\label{eq:pdfv2-terminal-evaluacion-ponderada}
 \widehat{\mathsf E}_{k,N}^{\mu}e_{\widehat x_k}
 :=\sum_{h\in\mathscr H_{k,N}(\widehat x_k)}
 \mu_{\widehat x_k,N}(h)^{1/2}
 e_{\mathfrak r_{k,N}(h)}.
\end{equation}
The projectivity gives
\begin{equation}
\label{eq:pdfv2-terminal-isometria}
 \left\|\widehat{\mathsf E}_{k,N}^{\mu}e_{\widehat x_k}\right\|^2
 =\sum_{h\in\mathscr H_{k,N}(\widehat x_k)}
   \mu_{\widehat x_k,N}(h)=1.
\end{equation}
The sum is not substituted by ``the history that generates the cylinder''.

\section{The nonadic connection as correspondence}

The nonadic holonomy was generated by the odometer and the TPK connection,
without a membership hypothesis, in
--
. In this chapter\eqref{eq:continuo-comun-gamma9-apice} we use
\eqref{eq:continuo-comun-gamma9-relacion}
\begin{equation}
\label{eq:pdfv2-terminal-apice-gamma9}
 \mathsf Z_{9,k}:=\mathsf Z^{\Gamma}_{9,k}.
\end{equation}
The maps
\begin{equation}
\label{eq:pdfv2-terminal-span-gamma9}
 \widehat{\mathcal X}^{\rm tot}_k
 \xleftarrow{\ s_{9,k}\ }
 \mathsf Z^{\Gamma}_{9,k}
 \xrightarrow{\ t_{9,k}\ }
 \widehat{\mathcal X}^{\rm tot}_{k+9}
\end{equation}
are those of \eqref{eq:continuo-comun-gamma9-span}; its source map is surjective by \cref{prop:continuo-comun-gamma9-total-natural}. Every point of the apex preserves the nine phase-compatible emissions. It is not replaced by the pair of its endpoints.

For \(N\geq k+9\), define the apex of horizon before of linearizing:
\begin{equation}
\label{eq:pdfv2-terminal-apice-horizonte-gamma9}
 \mathsf Z^\Gamma_{9;k,N}:=
 \coprod_{z=(\widehat x_k,\boldsymbol\omega)
            \in\mathsf Z^\Gamma_{9,k}}
 \operatorname{Term}_{k+9,N}(t_{9,k}z).
\end{equation}
An element is written \((z,q)\), where \(\boldsymbol\omega\) is the
nonadic word of nine emissions and \(q\) its tail up to \(N\). The maps on
histories use the labeled codomain
\begin{equation}
\label{eq:pdfv2-terminal-gamma9-codominio-colas}
 \mathscr D^\Gamma_{9;k,N}:=
 \coprod_{z\in\mathsf Z^\Gamma_{9,k}}
 \operatorname{Term}_{k+9,N}(t_{9,k}z).
\end{equation}
Thus the maps are
\begin{align}
 S_{k,N}(z,q)&:=\boldsymbol\omega\star q,
 \label{eq:pdfv2-terminal-gamma9-s-horizonte}\\
 T_{k,N}(z,q)&:=(z,q)\in\mathscr D^\Gamma_{9;k,N}.
 \label{eq:pdfv2-terminal-gamma9-t-horizonte}
\end{align}
The first lands in
\(\coprod_x\mathscr H_{k,N}(x)\); the second expressly preserves the
labelled summand by \(z\).

The link map of the apex is the explicit action
\begin{equation}
\label{eq:pdfv2-terminal-gamma9-enlace-apice}
 \pi^Z_{N+1,N}(z,q\star\varepsilon):=(z,q).
\end{equation}
It is surjective: if \((z,q)\) is given, one adds to the tail the emission
neutral at the endpoint of \(q\). The codomain of tails has link
\begin{equation}
\label{eq:pdfv2-terminal-gamma9-enlace-codominio}
 \pi^{\mathscr D}_{N+1,N}(z,q\star\varepsilon):=(z,q).
\end{equation}
We denote by
\(\pi^{\mathscr H,j}_{N+1,N}\) the map
\eqref{eq:pdfv2-terminal-truncamiento-historias} with initial level \(j\).
Then
\begin{align}
 \pi^{\mathscr H,k}_{N+1,N}S_{k,N+1}
 &=S_{k,N}\pi^Z_{N+1,N},
 \label{eq:pdfv2-terminal-gamma9-cuadrado-fuente}\\
 \pi^{\mathscr D}_{N+1,N}T_{k,N+1}
 &=T_{k,N}\pi^Z_{N+1,N}.
 \label{eq:pdfv2-terminal-gamma9-cuadrado-destino}
\end{align}
In the first square, both paths send
\(\boldsymbol\omega\star q\star\varepsilon\) to
\(\boldsymbol\omega\star q\); in the second, they send
\((z,q\star\varepsilon)\) to \((z,q)\). By unique decomposition of a word
into its first nine emissions and its tail, \(S_{k,N}\) is bijective. The
second square is Cartesian by the coproduct--fibered definition of
\eqref{eq:pdfv2-terminal-apice-horizonte-gamma9}. These are the
Beck--Chevalley squares preserving witnesses and multiplicities.

The corresponding apex of registers, now generated by the thirteen updates, is
\begin{equation}
\label{eq:pdfv2-terminal-apice-terminal-gamma9}
 \begin{split}
 \mathsf Z^{\mathfrak T}_{9;k,N}:=\bigl\{
 (&z,q;
   \mathfrak r_{k,N}(\boldsymbol\omega\star q),
   \mathfrak r_{k+9,N}(q)):\\
 &z=(\widehat x_k,\boldsymbol\omega)\in
       \mathsf Z^\Gamma_{9,k},\quad
 q\in\operatorname{Term}_{k+9,N}(t_{9,k}z)
 \bigr\}.
 \end{split}
\end{equation}
Its source and target maps read, respectively, the primer or the second
record, but never remove from their domain the witness \((z,q)\):
\begin{equation}
\label{eq:pdfv2-terminal-span-terminal}
 \mathfrak T_{k,N}
 \xleftarrow{\ s^{\mathfrak T}_{9;k,N}\ }
 \mathsf Z^{\mathfrak T}_{9;k,N}
 \xrightarrow{\ t^{\mathfrak T}_{9;k,N}\ }
 \mathfrak T_{k+9,N}.
\end{equation}
If \(\zeta=(z,q;r_s,r_t)\), its action is
\begin{equation}
\label{eq:pdfv2-terminal-span-terminal-accion}
 s^{\mathfrak T}_{9;k,N}(\zeta):=r_s,
 \qquad
 t^{\mathfrak T}_{9;k,N}(\zeta):=r_t.
\end{equation}
Its link is defined by
\begin{equation}
\label{eq:pdfv2-terminal-gamma9-enlace-registros}
 \begin{split}
 \pi^{Z,\mathfrak T}_{N+1,N}
 \bigl(&z,q\star\varepsilon;
       \mathfrak r_{k,N+1}
          (\boldsymbol\omega\star q\star\varepsilon),
       \mathfrak r_{k+9,N+1}(q\star\varepsilon)\bigr)\\
 &:=(z,q;\mathfrak r_{k,N}(\boldsymbol\omega\star q),
       \mathfrak r_{k+9,N}(q)).
 \end{split}
\end{equation}
The pair \((z,q)\) is a coordinate of the apex, not a representation
chosen of a triad; therefore the link is well defined even if two
registros terminals coincide. The compatibility of
\(\operatorname{Upd}\) with the deletion of the ultima edge proof that
\eqref{eq:pdfv2-terminal-gamma9-enlace-registros} commutes with both maps of
\eqref{eq:pdfv2-terminal-span-terminal}.

If an action on free modules is desired, this
correspondence is linearized by summing all witnesses:
\begin{equation}
\label{eq:pdfv2-terminal-linealizacion-span}
 [\Gamma_9]_{k,N}[r]
 :=\sum_{\zeta\in(s^{\mathfrak T}_{9;k,N})^{-1}(r)}
 [t^{\mathfrak T}_{9;k,N}(\zeta)].
\end{equation}
The formula preserves multiplicities. Equality of phases after
nine steps does not contract the states, the witnesses, or their ledgers.

\section{Uniqueness of the prefix decomposition and naturality of the truncation system}

For \(k\leq\ell\leq N\), every history of
\(\mathscr H_{k,N}(\widehat x_k)\) has a unique prefix
\(p\in\mathscr P_{\ell\mid k}(\widehat x_k)\). Therefore there is a disjoint decomposition
by prefixes, expressed without choosing any of them:
\begin{equation}
\label{eq:pdfv2-terminal-particion-historias}
 \mathscr H_{k,N}(\widehat x_k)
 =\bigsqcup_{p\in\mathscr P_{\ell\mid k}(\widehat x_k)}
 \left\{p\star q:q\in\mathscr H_{\ell,N}(p)\right\}.
\end{equation}
Let \(\operatorname{Conc}_p\) be the operation that prepends the record of the
prefix \(p\) and successively applies
\eqref{eq:pdfv2-terminal-actualizacion-unica}. We define the transition in common over a complete family of terminals by
\begin{equation}
\label{eq:pdfv2-terminal-transicion-comun}
 \mathfrak u_{\ell,k}
 \left(
   \bigl(\mathfrak T_{\ell,N}(p)\bigr)_{p\in\mathscr P_{\ell\mid k}(\widehat x_k)}
 \right)
 :=\bigcup_{p\in\mathscr P_{\ell\mid k}(\widehat x_k)}
   \operatorname{Conc}_p\mathfrak T_{\ell,N}(p).
\end{equation}

Writing
\(\mathbf G_{k,N}(\widehat x_k):=\mathfrak T_{k,N}(\widehat x_k)\), all naturality is concentrated in the single identity
\begin{equation}
\label{eq:pdfv2-terminal-naturalidad-unica}
 \boxed{
 \mathfrak u_{\ell,k}
 \left(
   \bigl(\mathbf G_{\ell,N}(p)\bigr)_{p\in\mathscr P_{\ell\mid k}(\widehat x_k)}
 \right)
 =\mathbf G_{k,N}(\widehat x_k).}
\end{equation}
There is no distinct law for each internal discriminant. The five
characteristics obey
\eqref{eq:pdfv2-terminal-naturalidad-unica} because each is updated
by means of the same prefix, the same extension, and the same ledger.

The nonadic correspondence is compatible with this identity through the
maps already constructed, not through a condition that is implicit. For every
horizon, \eqref{eq:pdfv2-terminal-gamma9-cuadrado-fuente} and
\eqref{eq:pdfv2-terminal-gamma9-cuadrado-destino} preserve the nonadic witness with its multiplicity, and
\eqref{eq:pdfv2-terminal-gamma9-enlace-registros} does the same in the
thirteen coordinates. If \(a\) is a coherent symmetry of the common
state, its action
\[
 a^Z(z,q):=(a^Zz,aq)
\]
satisfies
\begin{equation}
\label{eq:pdfv2-terminal-gamma9-equivarianza-completa}
 S a^Z=aS,\qquad T a^Z=aT,\qquad
 \pi^Za^Z=a^Z\pi^Z,\qquad
 s^{\mathfrak T}a^Z=as^{\mathfrak T},\qquad
 t^{\mathfrak T}a^Z=at^{\mathfrak T}.
\end{equation}
These equalities are verified entry by entry using
\eqref{eq:continuo-comun-gamma9-naturalidad} and the naturality of the thirteen
updates. They are the material condition that permits the use of
\eqref{eq:pdfv2-terminal-linealizacion-span}; a mere equality of visible
values does not replace it.

\section{Passage to the limit without interchanging image and limit}

Truncating a history and each of the thirteen fields defines
\begin{equation}
\label{eq:pdfv2-terminal-transicion-horizonte}
 \pi^{\mathfrak T}_{N+1,N}:
 \mathfrak T_{k,N+1}(\widehat x_k)\longrightarrow
 \mathfrak T_{k,N}(\widehat x_k),
 \qquad
 \pi^{\mathfrak T}_{N+1,N}\mathfrak r_{k,N+1}(h)
 =\mathfrak r_{k,N}(h_{\leq N}).
\end{equation}
It is surjective by
\eqref{eq:pdfv2-terminal-truncamiento-historias}. The complete terminal is
therefore
\begin{equation}
\label{eq:pdfv2-terminal-limite-inverso}
 \mathfrak T_{k,\infty}(\widehat x_k)
 :=\left\{(r_N)_{N\geq k}:
 r_N\in\mathfrak T_{k,N}(\widehat x_k),\quad
 \pi^{\mathfrak T}_{N+1,N}r_{N+1}=r_N\right\}.
\end{equation}
The finite sets of
\eqref{eq:pdfv2-terminal-multivaluado} are nonempty and their transition maps
are surjective; therefore,
\begin{equation}
\label{eq:pdfv2-terminal-limite-no-vacio}
 \mathfrak T_{k,\infty}(\widehat x_k)\neq\varnothing.
\end{equation}

\begin{corollary}[Simultaneous persistence of the five developments]
\label{cor:pdfv2-terminal-limite-cinco-desarrollos}
In each element of
\(\mathfrak T_{k,\infty}(\widehat x_k)\) there simultaneously persist the strong
limit \(\mathcal T_\infty\) of the polar telescopy, the Gram family and
losses, the incidence measure with its two naturalities, the cone
\(\operatorname{Cone}(I-U_{\Gamma_9,N})\), and the compatible family of Smith forms and
Bockstein pages. None of these data is obtained
by reconstructing only the visible terminal value.
\end{corollary}

\begin{proof}
Surjectivity of the truncations preserves the complete witnesses. The
telescoping, naturality, and reduction equations of
\cref{prop:pdfv2-terminal-cinco-desarrollos-integrales} commute with those
truncations; by universality of the inverse limit, they define a unique compatible
family on each surviving history.
\end{proof}

The apex nonadico in the limit is defined, not inferred from an image:
\begin{equation}
\label{eq:pdfv2-terminal-gamma9-apice-limite-raw}
 \mathsf Z^\Gamma_{9;k,\infty}
 :=\varprojlim_{N\geq k+9}
 \bigl(\mathsf Z^\Gamma_{9;k,N},\pi^Z_{N+1,N}\bigr).
\end{equation}
Since the links do not modify the finite witness \(z\) and only truncate its
tail, there exists the explicit isomorphism
\begin{equation}
\label{eq:pdfv2-terminal-gamma9-apice-limite-msect}
 \mathsf Z^\Gamma_{9;k,\infty}
 \xrightarrow{\ \cong\ }
 \coprod_{z\in\mathsf Z^\Gamma_{9,k}}
 \operatorname{Msect}(t_{9,k}z),\qquad
 ((z,q_N))_N\longmapsto(z,(q_N)_N).
\end{equation}
The inverse takes a witness \(z\) and a compatible tail \((q_N)_N\), and forms
\(((z,q_N))_N\). Both compositions are identities coordinate by
coordinate. The limit maps are
\begin{equation}
\label{eq:pdfv2-terminal-gamma9-mapas-limite-raw}
 S_{k,\infty}((z_N)_N):=(S_{k,N}z_N)_N,\qquad
 T_{k,\infty}((z_N)_N):=(T_{k,N}z_N)_N;
\end{equation}
are well defined precisely by \eqref{eq:pdfv2-terminal-gamma9-cuadrado-fuente} and \eqref{eq:pdfv2-terminal-gamma9-cuadrado-destino}.

The apex of registers is defined in the same way
\begin{equation}
\label{eq:pdfv2-terminal-gamma9-apice-registros-limite}
 \mathsf Z^{\mathfrak T}_{9;k,\infty}
 :=\varprojlim_{N\geq k+9}
 \bigl(\mathsf Z^{\mathfrak T}_{9;k,N},
       \pi^{Z,\mathfrak T}_{N+1,N}\bigr).
\end{equation}
The correspondences
\eqref{eq:pdfv2-terminal-span-terminal} then form, component by
component, the span
\begin{equation}
\label{eq:pdfv2-terminal-span-limite}
 \mathfrak T_{k,\infty}
 \xleftarrow{\ s^{\mathfrak T}_{9;k,\infty}\ }
 \mathsf Z^{\mathfrak T}_{9;k,\infty}
 \xrightarrow{\ t^{\mathfrak T}_{9;k,\infty}\ }
 \mathfrak T_{k+9,\infty}.
\end{equation}
The image of a limit has not been identified with the limit of certain
images: \eqref{eq:pdfv2-terminal-limite-inverso},
\eqref{eq:pdfv2-terminal-gamma9-apice-limite-raw} and
\eqref{eq:pdfv2-terminal-gamma9-apice-registros-limite} construct
directly the three compatible families and their maps.

\section{Internal reconstruction of the enriched state from its \(13\) genealogical coordinates}

Recovering does not unpack names: it repeats the fold that generated the
record. For
\(h=(\varepsilon_{k+1},\ldots,\varepsilon_N)\), we define
\begin{equation}
\label{eq:pdfv2-terminal-recuperacion-trece}
 \boxed{
 \operatorname{Rec}_{13;k,N}(h):=
 \operatorname{Upd}_{\varepsilon_N}\circ\cdots\circ
 \operatorname{Upd}_{\varepsilon_{k+1}}
 (\mathfrak r^0_{k,r}).}
\end{equation}
By \eqref{eq:pdfv2-terminal-registro-trece}, \(\operatorname{Rec}_{13;k,N}(h)=\mathfrak r_{k,N}(h)\), but the equality is now an identity between two compositions of actions, not the projection of a preserved input. If \(\operatorname{pr}_{\mathsf A}\) denotes the projector from the register to field \(\mathsf A\in\{\mathsf F,\mathsf E,\mathsf R,\mathsf N,\mathsf O,
\mathsf K,\mathsf M,\mathsf B,\mathsf G,\mathsf X,\mathsf S,\mathsf T,
\mathsf L\}\), then
\begin{equation}
\label{eq:pdfv2-terminal-recuperacion-campo-fold}
 \operatorname{pr}_{\mathsf A}\operatorname{Rec}_{13;k,N}(h)
 =\operatorname{Upd}^{\mathsf A}_{\varepsilon_N}\circ\cdots\circ
  \operatorname{Upd}^{\mathsf A}_{\varepsilon_{k+1}}
  (\mathsf A_{k,k}).
\end{equation}
The proof is induction on \(N-k\): for \(N=k\) both sides are the
seed; the inductive step applies the component \(\mathsf A\) of
\(\operatorname{Upd}_{\varepsilon_N}\).

Adem\'as,
\begin{equation}
\label{eq:pdfv2-terminal-recuperacion-compatible}
 \operatorname{Rec}_{13;k,N}(h_{\leq N})
 =\operatorname{tr}_{13,N+1,N}
   \bigl(\operatorname{Rec}_{13;k,N+1}(h)\bigr),
\end{equation}
where \(\operatorname{tr}_{13,N+1,N}\) removes only the last
update and restricts its cone of tails. Consequently, an element
of \(\mathfrak T_{k,\infty}(\widehat x_k)\) determines a compatible family of the
thirteen coordinates at every depth.

\section{Multivalued terminality, naturality, and compatible limit of the continuum}

\begin{theorem}[Multivalued terminal, naturality, and limit of the continuum]
\label{thm:pdfv2-terminal-conjunto}
For every \(\widehat x_k\in\widehat{\mathcal X}^{\rm tot}_k\), the rules
\eqref{eq:pdfv2-terminal-historias-finita}--
\eqref{eq:pdfv2-terminal-recuperacion-compatible} construct, from the
unique discrete continuum:
\begin{enumerate}
 \item a multisection terminal finite not empty at each horizon;
 \item an evaluation that conserves all histories and multiplicities;
 \item the connection nonadic as a correspondence with witnesses;
 \item the unique identity of naturality
       \eqref{eq:pdfv2-terminal-naturalidad-unica};
 \item a nonempty inverse terminal and a nonadic correspondence in the
       limit;
 \item the compatible recovering of the fiber, extensions, relations,
       numbers, orientation, carry, memory, boundary, route,
       multisection, survival, terminal, and reader.
\end{enumerate}
The five intrinsic characteristics remain correlated within
each record and each extension witness.
\end{theorem}

\begin{proof}
The non-emptiness of
\(\mathscr H_{k,N}(\widehat x_k)\) is nonempty by
\cref{prop:continuo-comun-arbol-finito-no-vacio}, and its truncation is
surjective by \cref{lem:continuo-comun-terminal-sobreyectivo}. The rule
\eqref{eq:pdfv2-terminal-actualizacion-unica} applies the APP,
TRIT, and TPK identities to each witness \(\varepsilon_j\); by induction on
\(N-k\), it produces all the fields of
\eqref{eq:pdfv2-terminal-registro-trece} and conserves their relations.

Taking the complete set in
\eqref{eq:pdfv2-terminal-multivaluado}, or the complete sum in
\eqref{eq:pdfv2-terminal-evaluacion-todas}, proves that the terminal does not
contain a selector. The weighted version is isometric by the projective
equality \eqref{eq:pdfv2-terminal-isometria}.

The witness \(z=(\widehat x_k,\boldsymbol\omega)\) carries the phase action of each of its nine emissions; by \cref{prop:continuo-comun-gamma9-retorno-memoria} it defines the two maps in \eqref{eq:pdfv2-terminal-span-gamma9} without converting the relation into a function or identifying phase with state. Equations \eqref{eq:pdfv2-terminal-gamma9-s-horizonte}--\eqref{eq:pdfv2-terminal-gamma9-enlace-registros} prove the two horizon squares and the same construction over registers. Summing the source space of each witness yields \eqref{eq:pdfv2-terminal-linealizacion-span} without erasing multiplicities.

The partition
\eqref{eq:pdfv2-terminal-particion-historias} is disjoint because every
history has a determined prefix of length \(\ell\). Concatenating each
prefix with all its continuations reconstructs once and only once each
history of \(\mathscr H_{k,N}(\widehat x_k)\). This directly proves
\eqref{eq:pdfv2-terminal-naturalidad-unica} for the whole register; five disconnected proofs are not
needed.

Surjectivity of
\eqref{eq:pdfv2-terminal-transicion-horizonte} and finiteness of each level
imply nonemptiness of
\eqref{eq:pdfv2-terminal-limite-inverso}. The explicit links
\eqref{eq:pdfv2-terminal-gamma9-enlace-apice} and
\eqref{eq:pdfv2-terminal-gamma9-enlace-registros}, together with their squares,
construct the limits
\eqref{eq:pdfv2-terminal-gamma9-apice-limite-raw} and
\eqref{eq:pdfv2-terminal-gamma9-apice-registros-limite}; from them follows
\eqref{eq:pdfv2-terminal-span-limite}. Finally,
\eqref{eq:pdfv2-terminal-recuperacion-trece} reruns the
successive updates, and
\eqref{eq:pdfv2-terminal-recuperacion-compatible} proves that those
readings pass jointly to the limit.
\end{proof}

\section{Obstructions to the terminality of the enriched state: loss of fiber, orientation, memory, or reader}

\begin{proposition}[Falsifiers of the terminal set]
\label{prop:pdfv2-terminal-falsadores}
Any of the following operations replaces the constructed object and
blocks the conclusion of
\cref{thm:pdfv2-terminal-conjunto}:
\begin{enumerate}
 \item send to cylinder to to single history when it admits several;
 \item choosing a tail of the terminal field \(\mathsf T\) and eliminating the
       remaining ones;
 \item replace \(\Gamma_9\) by a function without proving uniqueness of
       each target and of its witness;
 \item linearize to correspondence while ignoring multiplicities;
 \item identifying phase return with state return;
 \item preserve the input intact as the first coordinate and use its
       recoverability as purported generation;
 \item omit any of the thirteen fields of
       \eqref{eq:pdfv2-terminal-registro-trece};
 \item allowing two intrinsic characteristics to use different histories,
       carries, boundaries, or terminals;
 \item replace
       \eqref{eq:pdfv2-terminal-naturalidad-unica} by five independent
       assertions;
 \item assert operatorial compatibility of \(\Gamma_9\) without preserving the
       space of witnesses and its multiplicities;
 \item interchanging, without proof, formation of images and passage to the limit;
 \item present a list of thirteen names without the actions
       \eqref{eq:pdfv2-terminal-actualizacion-unica} that produce them.
 \item omit the polar telescoping, the loss row, the naturality
       adjoint, the module of all continuations or the reduction
       Smith--Bockstein certified in
       \cref{prop:pdfv2-terminal-cinco-desarrollos-integrales};
 \item turning any of those five developments into an autonomous branch
       or using a later classical application to select its domain.
\end{enumerate}
\end{proposition}

\begin{proof}
The first five operations contract the multisection or the correspondence. The
next four break their common support and naturality. The final
three eliminate the generative proof or alter passage to the limit. In
each case, the resulting object ceases to satisfy at least one of
\eqref{eq:pdfv2-terminal-multivaluado},
\eqref{eq:pdfv2-terminal-span-terminal},
\eqref{eq:pdfv2-terminal-naturalidad-unica}, or
\eqref{eq:pdfv2-terminal-recuperacion-compatible}; therefore, the conclusion
does not carry over.
\end{proof}

\begin{statusbox}
The organization of the complete terminal object, unique naturality, and passage to the
limit are established jointly. The surviving multisection, the
carries, the register, the boundary, the route, and the nonadic connection are its
structural antecedents.\quad
The linearization over free modules is proved by
\eqref{eq:pdfv2-terminal-linealizacion-span} and its Beck--Chevalley
squares. A subsequent promotion to a bounded operator on an
analytic completion additionally requires that norm to be declared and the bound proved; the
correspondence \eqref{eq:pdfv2-terminal-span-limite} does not depend on that
interface.
\end{statusbox}
